\documentclass[11pt]{article}
\usepackage[margin=1in]{geometry}
\usepackage{xcolor}
\usepackage[breakable,listings]{tcolorbox}
\tcbuselibrary{listings}
\usepackage{hyperref}
\usepackage{enumitem}
\usepackage{listings}
\usepackage{verbatim}
\usepackage{amsmath}
\usepackage{graphicx}
\usepackage{booktabs}
\usepackage{float}

% Define colors
\definecolor{osaorange}{RGB}{220,115,38}
\definecolor{codebackground}{RGB}{248,248,248}

% Configure hyperref
\hypersetup{
    colorlinks=true,
    linkcolor=blue,
    urlcolor=blue,
    citecolor=blue
}

% Configure listings for code display
\lstset{
  basicstyle=\small\ttfamily,
  breaklines=true,
  breakatwhitespace=true,
  postbreak=\mbox{\textcolor{gray}{$\hookrightarrow$}\space},
  columns=flexible,
  keepspaces=true,
  showstringspaces=false
}

% Code environment using tcolorbox with listings
\newtcblisting{rcode}{
  colback=white,
  colframe=gray!50,
  boxrule=0.5pt,
  arc=0pt,
  left=3pt,
  right=3pt,
  top=3pt,
  bottom=3pt,
  width=\textwidth,
  breakable,
  listing only,
  listing options={
    basicstyle=\small\ttfamily,
    breaklines=true,
    breakatwhitespace=true,
    postbreak=\mbox{\textcolor{gray}{$\hookrightarrow$}\space},
    columns=flexible,
    keepspaces=true,
    showstringspaces=false
  }
}

\title{}
\author{}
\date{}

\begin{document}

% Header box
\begin{tcolorbox}[colback=osaorange,colframe=black,boxrule=2pt,arc=0pt,left=3pt,right=3pt,top=6pt,bottom=6pt,boxsep=2pt]
\centering
\textcolor{white}{\textbf{\Large H524 Introduction to Biostatistics}}\\
\textcolor{white}{\textbf{\Large Week 3 Lab: Probability Distributions and Sampling}}\\
\textcolor{white}{\textbf{\Large Fall 2025}}
\end{tcolorbox}

\vspace{8pt}

\noindent\textbf{Format:} Online\\
\textbf{Tools Required:} R, RStudio\\
\textbf{Estimated Time:} 90 minutes

\section{Learning Objectives}

By the end of this lab, you will be able to:
\begin{enumerate}
\item Work with the binomial distribution in R (\texttt{dbinom}, \texttt{pbinom})
\item Calculate probabilities using the normal distribution (\texttt{dnorm}, \texttt{pnorm}, \texttt{qnorm})
\item Understand standardization (Z-scores) and use R to convert between scales
\item Simulate sampling distributions to demonstrate the Central Limit Theorem
\item Calculate and interpret confidence intervals in R
\item Visualize probability distributions and sampling distributions
\end{enumerate}

\section{Part 1: Binomial Distribution in R}

The binomial distribution models the number of successes in a fixed number of independent trials, each with the same probability of success.

\subsection{Exercise 1.1: Basic Binomial Calculations}

\textbf{Scenario:} A vaccine has 85\% efficacy. We vaccinate 20 people.

\begin{rcode}
# Binomial distribution: n=20, p=0.85
n <- 20
p <- 0.85

# Probability of exactly 17 protected
prob_17 <- dbinom(17, size=n, prob=p)
cat("P(X = 17):", round(prob_17, 4), "\n")

# Probability of 17 or fewer protected
prob_at_most_17 <- pbinom(17, n, p)
cat("P(X <= 17):", round(prob_at_most_17, 4), "\n")

# Probability of at least 18 protected
prob_at_least_18 <- pbinom(17, n, p, lower.tail=FALSE)
# OR: 1 - pbinom(17, n, p)
cat("P(X >= 18):", round(prob_at_least_18, 4), "\n")

# Expected value and standard deviation
expected <- n * p
std_dev <- sqrt(n * p * (1 - p))
cat("\nExpected value:", expected, "\n")
cat("Standard deviation:", round(std_dev, 3), "\n")
\end{rcode}

\textbf{Try it yourself:} What if the efficacy was 90\% instead? Recalculate all the above probabilities with \texttt{p = 0.90}.

\vspace{2cm}

\subsection{Exercise 1.2: Visualizing the Binomial Distribution}

\begin{rcode}
# Create a visualization of the binomial distribution
n <- 20
p <- 0.85

# Create vector of possible values
x_values <- 0:n

# Calculate probabilities for each value
probabilities <- dbinom(x_values, n, p)

# Create barplot and capture bar positions
bp <- barplot(probabilities,
              names.arg = x_values,
              main = "Binomial Distribution (n=20, p=0.85)",
              xlab = "Number of Successes",
              ylab = "Probability",
              col = "lightblue",
              border = "white")

# Add vertical line at expected value
abline(v = bp[which(x_values == n*p)], col = "red", lwd = 2, lty = 2)

# Add legend
legend("topleft",
       legend = "Expected Value",
       col = "red",
       lty = 2,
       lwd = 2)
\end{rcode}

\textbf{Question:} Look at your plot. Does it appear symmetric or skewed? Why?

\vspace{1.5cm}

\subsection{Exercise 1.3: Binomial Scenarios}

\begin{rcode}
# Antibiotic treatment: 75% success rate, treat 15 patients
antibiotic_prob <- dbinom(12, size=15, prob=0.75)
cat("P(exactly 12 successes):", round(antibiotic_prob, 4), "\n")

# At least 10 successes
at_least_10 <- sum(dbinom(10:15, size=15, prob=0.75))
# OR: pbinom(9, 15, 0.75, lower.tail=FALSE)
cat("P(at least 10 successes):", round(at_least_10, 4), "\n")
\end{rcode}

\textbf{Practice:} A clinical trial tests a drug on 25 patients. The success rate is 0.70. Find:
\begin{enumerate}[label=\alph*)]
\item P(exactly 18 successes)
\item P(fewer than 15 successes)
\item P(between 15 and 20 successes, inclusive)
\end{enumerate}

\vspace{2cm}

\section{Part 2: Normal Distribution in R}

The normal distribution is the most important continuous distribution in statistics. R provides four functions: \texttt{dnorm} (density), \texttt{pnorm} (cumulative probability), \texttt{qnorm} (quantiles), and \texttt{rnorm} (random generation).

\subsection{Exercise 2.1: Basic Normal Calculations}

\textbf{Scenario:} Cholesterol levels are normally distributed with mean 200 mg/dL and SD 40 mg/dL.

\begin{rcode}
# Normal distribution parameters
mu <- 200      # mean
sigma <- 40    # standard deviation

# Probability that cholesterol < 240
prob_below_240 <- pnorm(240, mean=mu, sd=sigma)
cat("P(X < 240):", round(prob_below_240, 4), "\n")

# Probability that cholesterol > 240
prob_above_240 <- pnorm(240, mu, sigma, lower.tail=FALSE)
# OR: 1 - pnorm(240, mu, sigma)
cat("P(X > 240):", round(prob_above_240, 4), "\n")

# Probability between 180 and 220
prob_between <- pnorm(220, mu, sigma) - pnorm(180, mu, sigma)
cat("P(180 < X < 220):", round(prob_between, 4), "\n")
\end{rcode}

\textbf{Question:} About what proportion of people have cholesterol above 240 mg/dL (considered high)?

\vspace{1cm}

\subsection{Exercise 2.2: Finding Percentiles}

\begin{rcode}
# Find the 90th percentile of cholesterol
percentile_90 <- qnorm(0.90, mean=mu, sd=sigma)
cat("90th percentile:", round(percentile_90, 1), "mg/dL\n")

# Find the 25th, 50th (median), and 75th percentiles
quartiles <- qnorm(c(0.25, 0.50, 0.75), mu, sigma)
names(quartiles) <- c("Q1", "Median", "Q3")
print(round(quartiles, 1))

# Verify: what proportion is below the 90th percentile?
check <- pnorm(percentile_90, mu, sigma)
cat("\nVerification:", round(check, 2), "\n")
\end{rcode}

\textbf{Try it yourself:} Find the cholesterol levels that define:
\begin{itemize}
\item The bottom 10\% (10th percentile)
\item The top 5\% (95th percentile)
\end{itemize}

\vspace{1.5cm}

\subsection{Exercise 2.3: Standardization (Z-scores)}

\begin{rcode}
# Convert a cholesterol value to a Z-score
chol_value <- 260   # mg/dL

# Calculate Z-score manually
z_score <- (chol_value - mu) / sigma
cat("Z-score for", chol_value, "mg/dL:", round(z_score, 2), "\n")

# Interpretation
cat("This value is", round(z_score, 2),
    "standard deviations above the mean\n")

# Use the standard normal to find probability
prob <- pnorm(z_score)
cat("Proportion below", chol_value, ":", round(prob, 4), "\n")

# Alternative: use original scale directly
prob_direct <- pnorm(chol_value, mu, sigma)
cat("Direct calculation:", round(prob_direct, 4), "\n")

# R also has a scale() function for vectors
cholesterol_data <- c(180, 220, 195, 240, 175, 210)
z_scores <- scale(cholesterol_data, center=mu, scale=sigma)
print(round(z_scores, 2))
\end{rcode}

\textbf{Question:} If someone has a Z-score of 1.5, what is their actual cholesterol level?

\vspace{1.5cm}

\subsection{Exercise 2.4: Visualizing the Normal Distribution}

\begin{rcode}
# Create a normal distribution plot
x <- seq(100, 300, length=200)
y <- dnorm(x, mean=mu, sd=sigma)

plot(x, y, type="l", lwd=2, col="blue",
     main="Normal Distribution of Cholesterol",
     xlab="Cholesterol (mg/dL)",
     ylab="Density")

# Add vertical lines for mean and +/- 1 SD
abline(v=mu, col="red", lwd=2, lty=1)
abline(v=c(mu-sigma, mu+sigma), col="red", lwd=1, lty=2)

# Shade region above 240
x_high <- seq(240, 300, length=100)
y_high <- dnorm(x_high, mu, sigma)
polygon(c(240, x_high, 300), c(0, y_high, 0),
        col=rgb(1,0,0,0.3), border=NA)

# Add text
text(mu, max(y)*0.9, "Mean", pos=3)
text(240, max(y)*0.5, "High\nCholesterol", pos=4)

# Add legend
legend("topright",
       legend=c("Mean", "+/- 1 SD"),
       col="red",
       lty=c(1,2),
       lwd=c(2,1))
\end{rcode}

\section{Part 3: Sampling Distributions and Central Limit Theorem}

Now we'll see one of the most important concepts in statistics: the Central Limit Theorem (CLT).

\subsection{Exercise 3.1: Understanding Standard Error}

\begin{rcode}
# Population parameters
pop_mean <- 120        # mean blood pressure
pop_sd <- 20           # population SD

# Calculate standard error for different sample sizes
sample_sizes <- c(5, 10, 25, 50, 100)

for (n in sample_sizes) {
  se <- pop_sd / sqrt(n)
  cat("n =", n, "\t SE =", round(se, 2), "\n")
}

# Notice: SE decreases as n increases!
\end{rcode}

\textbf{Question:} To cut the standard error in half, by what factor must you increase the sample size?

\vspace{1.5cm}

\subsection{Exercise 3.2: Simulating the Sampling Distribution}

\begin{rcode}
# Simulate drawing many samples and computing their means
set.seed(524)

pop_mean <- 120
pop_sd <- 20
sample_size <- 30
num_samples <- 1000

# Draw 1000 samples, each of size 30, and compute their means
sample_means <- replicate(num_samples, {
  sample_data <- rnorm(sample_size, mean=pop_mean, sd=pop_sd)
  mean(sample_data)
})

# Summary of the sampling distribution
cat("Mean of sample means:", round(mean(sample_means), 2), "\n")
cat("SD of sample means (empirical SE):", round(sd(sample_means), 2), "\n")
cat("Theoretical SE:", round(pop_sd/sqrt(sample_size), 2), "\n")

# Visualize the sampling distribution
hist(sample_means, breaks=30, prob=TRUE,
     main=paste("Sampling Distribution of Mean (n=", sample_size, ")", sep=""),
     xlab="Sample Mean",
     ylab="Density",
     col="lightblue",
     border="white")

# Overlay theoretical normal distribution
se_theoretical <- pop_sd / sqrt(sample_size)
curve(dnorm(x, mean=pop_mean, sd=se_theoretical),
      add=TRUE, col="red", lwd=2)

legend("topright",
       legend="Theoretical (CLT)",
       col="red",
       lwd=2)
\end{rcode}

\textbf{Observation:} The distribution of sample means is approximately normal, even though we're sampling from a normal population. The CLT guarantees this!

\subsection{Exercise 3.3: CLT with Non-Normal Data (Exponential Population)}

Let's see the power of the CLT: it works even when the population is NOT normal!

\textbf{Step 1: Define a population that is exponential (clearly NOT normal)}

\begin{rcode}
# Define our population: exponential distribution
# This represents waiting times, highly right-skewed
set.seed(123)

# Population parameters
pop_mean <- 50        # mean of exponential
rate_param <- 1/50    # rate parameter for rexp()

# Show what the population looks like
population <- rexp(10000, rate=rate_param)

# Plot the population distribution
hist(population, breaks=50, prob=TRUE,
     main="Population Distribution (Exponential)",
     xlab="Value",
     col="lightcoral",
     border="white")

# Add population density curve
curve(dexp(x, rate=rate_param), add=TRUE, col="darkred", lwd=2)

# Add statistics
text(100, 0.015, paste("Mean =", round(mean(population), 1)), pos=4)
text(100, 0.013, paste("SD =", round(sd(population), 1)), pos=4)
text(100, 0.011, "Clearly SKEWED!", pos=4, font=2)
\end{rcode}

\textbf{Observation:} The population is highly right-skewed (NOT normal at all!)

\vspace{0.5cm}

\textbf{Step 2: Repeatedly sample from this population and calculate means}

\begin{rcode}
# Now let's see what happens when we take samples and compute means
set.seed(123)

# We'll try different sample sizes
sample_sizes <- c(5, 10, 30, 100)

# Compare sampling distributions for different sample sizes
par(mfrow=c(2,2))  # 2x2 plot layout

for (n in sample_sizes) {
  # Take 1000 samples of size n
  # For each sample, calculate the mean
  sample_means <- replicate(1000, {
    sample_data <- sample(population, n, replace=TRUE)
    mean(sample_data)
  })

  # Plot the distribution of sample means
  hist(sample_means, breaks=30, prob=TRUE,
       main=paste("Sample size n =", n),
       xlab="Sample Mean",
       col="lightblue",
       border="white",
       xlim=c(0, 100))

  # Overlay the theoretical normal distribution (CLT prediction)
  se_theoretical <- pop_mean / sqrt(n)  # SE for exponential
  curve(dnorm(x, mean=pop_mean, sd=se_theoretical),
        add=TRUE, col="red", lwd=2)

  # Add text
  text(70, max(hist(sample_means, plot=FALSE)$density) * 0.8,
       paste("SE =", round(se_theoretical, 2)),
       col="red")
}

par(mfrow=c(1,1))  # Reset plot layout
\end{rcode}

\textbf{Step 3: Observe that sample means ARE normally distributed!}

\textbf{Amazing CLT result:}
\begin{itemize}
\item \textbf{Population:} Highly skewed exponential (definitely NOT normal)
\item \textbf{Sample means:} Approximately normal distribution!
\item \textbf{As $n$ increases:} The approximation gets better and better
\item Even at $n=30$: Already looks quite normal
\item At $n=100$: Almost perfectly normal!
\end{itemize}

\vspace{0.5cm}

\textbf{Key Insight:} The CLT tells us that \textit{no matter what the population distribution looks like}, sample means will be approximately normally distributed for large enough $n$. This is why the normal distribution is so fundamental to statistical inference!

\subsection{Exercise 3.4: Probability Calculations with Sample Means}

\begin{rcode}
# Population: cholesterol with mean=200, SD=40
# Sample size: n=64

pop_mean <- 200
pop_sd <- 40
n <- 64

# Standard error
se <- pop_sd / sqrt(n)
cat("Standard Error:", se, "\n")

# Probability that sample mean > 210
prob_above_210 <- pnorm(210, mean=pop_mean, sd=se, lower.tail=FALSE)
cat("P(sample mean > 210):", round(prob_above_210, 4), "\n")

# Probability that sample mean is between 195 and 205
prob_between <- pnorm(205, pop_mean, se) - pnorm(195, pop_mean, se)
cat("P(195 < sample mean < 205):", round(prob_between, 4), "\n")

# Compare with individual probability
prob_individual <- pnorm(205, pop_mean, pop_sd) - pnorm(195, pop_mean, pop_sd)
cat("\nFor comparison, P(195 < individual < 205):",
    round(prob_individual, 4), "\n")
cat("Sample means are MUCH more concentrated around the population mean!\n")
\end{rcode}

\section{Part 4: Confidence Intervals in R}

Confidence intervals quantify our uncertainty about population parameters.

\subsection{Exercise 4.1: Manual CI Calculation}

\begin{rcode}
# Sample data: tumor growth measurements (mm)
tumor_growth <- c(7, 10, 9, 8, 7, 6, 8, 9, 12, 13)

# Calculate sample statistics
xbar <- mean(tumor_growth)
s <- sd(tumor_growth)
n <- length(tumor_growth)

cat("Sample mean:", round(xbar, 2), "\n")
cat("Sample SD:", round(s, 3), "\n")
cat("Sample size:", n, "\n\n")

# Calculate standard error
se <- s / sqrt(n)
cat("Standard error:", round(se, 3), "\n\n")

# 95% Confidence Interval using t-distribution
alpha <- 0.05
df <- n - 1
t_critical <- qt(1 - alpha/2, df)

cat("Degrees of freedom:", df, "\n")
cat("t critical value:", round(t_critical, 3), "\n\n")

# Margin of error
margin_error <- t_critical * se

# Confidence interval
ci_lower <- xbar - margin_error
ci_upper <- xbar + margin_error

cat("95% Confidence Interval:\n")
cat("(", round(ci_lower, 2), ", ", round(ci_upper, 2), ")\n", sep="")
\end{rcode}

\textbf{Interpretation:} We are 95\% confident that the true mean tumor growth is between 7.30 and 10.50 mm.

\subsection{Exercise 4.2: Using t.test() for CIs}

\begin{rcode}
# Much easier method: use t.test()
result <- t.test(tumor_growth, conf.level=0.95)

# Extract confidence interval
ci <- result$conf.int
cat("95% CI from t.test():", round(ci, 2), "\n")

# The t.test() function gives us additional info
print(result)

# For 99% CI, just change conf.level
result_99 <- t.test(tumor_growth, conf.level=0.99)
cat("\n99% CI:", round(result_99$conf.int, 2), "\n")
\end{rcode}

\textbf{Note:} The \texttt{t.test()} function is the easiest way to get confidence intervals in R!

\subsection{Exercise 4.3: Comparing Different Confidence Levels}

\begin{rcode}
# Compare 90%, 95%, and 99% confidence intervals
confidence_levels <- c(0.90, 0.95, 0.99)

for (conf in confidence_levels) {
  ci <- t.test(tumor_growth, conf.level=conf)$conf.int
  width <- ci[2] - ci[1]

  cat(conf*100, "% CI: (", round(ci[1], 2), ", ", round(ci[2], 2),
      ")  Width:", round(width, 2), "\n", sep="")
}

cat("\nNotice: Higher confidence = wider interval\n")
\end{rcode}

\subsection{Exercise 4.4: Effect of Sample Size on CIs}

\begin{rcode}
# Generate larger sample from same population
set.seed(42)

# Small sample (n=10)
small_sample <- rnorm(10, mean=100, sd=15)
ci_small <- t.test(small_sample)$conf.int
width_small <- ci_small[2] - ci_small[1]

# Large sample (n=100)
large_sample <- rnorm(100, mean=100, sd=15)
ci_large <- t.test(large_sample)$conf.int
width_large <- ci_large[2] - ci_large[1]

cat("Small sample (n=10):\n")
cat("  95% CI:", round(ci_small, 2), "\n")
cat("  Width:", round(width_small, 2), "\n\n")

cat("Large sample (n=100):\n")
cat("  95% CI:", round(ci_large, 2), "\n")
cat("  Width:", round(width_large, 2), "\n\n")

cat("Larger sample gives narrower (more precise) CI!\n")
cat("Width ratio:", round(width_small/width_large, 2), "\n")
\end{rcode}

\section{Practice Problems}

\subsection{Problem 1: Vaccine Efficacy}

A COVID-19 vaccine has 92\% efficacy. A clinic vaccinates 50 people.

\begin{enumerate}[label=\alph*)]
\item What is the probability that at least 45 people are protected?
\item What is the expected number protected?
\item Create a visualization of this binomial distribution
\end{enumerate}

\vspace{2cm}

\subsection{Problem 2: Birth Weights}

Birth weights are normally distributed with mean 3400g and SD 500g.

\begin{enumerate}[label=\alph*)]
\item What proportion of babies are born with low birth weight ($<$ 2500g)?
\item What birth weight represents the 10th percentile?
\item What proportion of babies have birth weights between 3000g and 4000g?
\end{enumerate}

\vspace{2cm}

\subsection{Problem 3: Sampling Distribution}

A population has mean blood glucose of 100 mg/dL with SD 20 mg/dL. You take a sample of 36 patients.

\begin{enumerate}[label=\alph*)]
\item What is the distribution of the sample mean?
\item What is the probability the sample mean exceeds 105 mg/dL?
\item Within what range will 95\% of sample means fall?
\end{enumerate}

\vspace{2cm}

\subsection{Problem 4: Confidence Interval}

A study measures pain reduction scores for 25 patients: mean reduction = 5.2 points, SD = 2.1 points.

\textbf{Note:} Use the given summary statistics directly. Calculate the CI manually using the formulas from Exercise 4.1.

\begin{enumerate}[label=\alph*)]
\item Calculate a 95\% confidence interval for the true mean reduction
\item Interpret this interval in one complete sentence
\item Would a 90\% CI be wider or narrower? By how much?
\end{enumerate}

\vspace{2cm}

\subsection{Problem 5: CLT Simulation}

Simulate the Central Limit Theorem:

\begin{enumerate}[label=\alph*)]
\item Generate 1000 samples of size 40 from a population with mean=75 and SD=12
\item Plot the histogram of sample means
\item Compare the empirical SD of sample means to the theoretical SE
\item Overlay the theoretical normal curve predicted by CLT
\end{enumerate}

\vspace{2cm}

\section{Saving Your Work}

\textbf{IMPORTANT:} Save your R script file!

\begin{rcode}
# At the top of your R script, add:
# Name: Your Name
# Date: Today's Date
# Lab: Week 3 - Probability Distributions and Sampling

# Save your workspace
save.image("week3_lab.RData")

# To load it later:
# load("week3_lab.RData")
\end{rcode}

\end{document}
